\chapter{Related Work}\label{cap:related_work}

This chapter reviews the literature most relevant to the development of automated methods for image-based wildlife monitoring, with particular emphasis on camera trap analysis, computer vision, deep learning, and multimodal learning. Its purpose is to situate the research problem addressed in this thesis within the broader methodological landscape, identifying the principal advances that have enabled large-scale ecological image interpretation as well as the limitations that remain unresolved. The reviewed works cover a range of complementary directions, including manual and semi-automated camera trap workflows, supervised species recognition, object detection pipelines, hierarchical and fine-grained classification strategies, metadata-enhanced models, zero-shot and vision--language approaches, and operational systems for ecological monitoring. Examining these areas is essential because the effectiveness of wildlife recognition systems depends not only on model architecture, but also on issues such as annotation cost, class imbalance, domain shift, contextual information, and deployment constraints in real field conditions. To provide a coherent account of this landscape, the chapter is organized into thematic sections that progressively examine these research directions and highlight the main challenges that motivate the approach proposed in this thesis.


\section{Manual Interpretation and Sources of Variability in Camera Trap Workflows}

Before discussing automated approaches, it is important to note that manual interpretation itself introduces variability into ecological monitoring workflows. \citep{zett2022inter} examine the reliability of wildlife information extracted from camera trap images by multiple human observers, addressing an often overlooked source of variability in ecological monitoring workflows. The study evaluates inter-observer agreement in the manual interpretation of camera trap images, focusing on species detection, species identification, and estimates of herbivore group sizes and individual counts. Ten observers with varying levels of experience independently processed a dataset of 11,560 camera trap images collected at seven waterhole locations in Ongava Game Reserve, Namibia, containing records of 22 mammal species. The analysis quantified differences in image processing speed, detection rates, and classification accuracy using statistical measures such as intra-class correlation coefficients. Results showed substantial variation among observers, with species detection rates ranging from 81.8\% to 100\% and overall correct detections reaching approximately 96\% of possible records. Misidentifications were most frequent among phenotypically similar species, and counting estimates for social herbivores exhibited high variance across observers. The study highlights that observer experience significantly influences accuracy and processing speed, and emphasizes that human interpretation can introduce bias in camera trap datasets, which has implications for biodiversity assessments and population monitoring workflows \citep{zett2022inter}.

\section{Supervised Deep Learning for Camera Trap Species Recognition}

A substantial body of work has addressed wildlife recognition in camera trap images through supervised deep learning, particularly with convolutional neural networks. \citep{binta2023animal} propose a deep learning framework for automated species recognition in camera trap images, addressing the challenge of manually processing large volumes of wildlife monitoring data. The study focuses on classifying herpetofauna—specifically snakes, lizards, and frogs/toads—from camera trap images collected across multiple locations in Texas, USA. The authors evaluate three convolutional neural network architectures: a custom CNN trained from scratch and two transfer learning models based on VGG16 and ResNet50. The dataset contains several thousand annotated images captured under varying environmental conditions, including day and night scenes, occlusions, and small object sizes. To address class imbalance and limited data, the authors apply oversampling and image augmentation during preprocessing. The results show that transfer learning improves classification performance, with VGG16 achieving approximately 87\% accuracy and ResNet50 86\%, both outperforming the custom CNN model. The study shows the potential of pretrained convolutional networks for species recognition in camera trap images, although the approach is limited by the relatively small and imbalanced dataset and by its focus on broad taxonomic groups rather than fine-grained species classification \citep{binta2023animal}.

Extending this line of research to a broader taxonomic scope, \citep{schneider2024recognition} propose a deep learning framework for automated recognition of wildlife species in camera trap images, addressing the limitation of existing systems that typically focus only on mammals or treat birds as a single category. Their approach integrates an object detection stage using MegaDetector to localize animals, followed by image classification models based on several modern neural architectures, including ResNet, EfficientNetV2, Vision Transformer, Swin Transformer, and ConvNeXt. The models are trained to recognize 25 mammal and 63 bird species commonly found in Central European forests and additionally predict hierarchical taxonomic labels such as genus, family, and order to exploit partially annotated datasets. The training data were aggregated from multiple public datasets and web sources, while evaluation was conducted on camera trap images collected in the Marburg Open Forest (Germany) and Białowieża National Park (Poland). The best-performing ConvNeXt model achieved a mean average precision of 97.91\% on validation data and 90.39\% and 82.77\% on the two test sets. The study demonstrates the feasibility of joint mammal and bird recognition in camera trap images, although performance decreases under domain shifts such as lower image quality or different camera placements \citep{schneider2024recognition}.

In addition to multi-species recognition, several studies focus on single-species detection as a means to analyze methodological choices in challenging ecological scenarios. \citep{burkard2024automated} investigate automated single-species detection in camera trap images using deep learning, focusing on the identification of grey herons (Ardea cinerea) in a challenging ecological monitoring dataset. The study analyzes how architectural choices, data handling strategies, and evaluation metrics affect model performance in the presence of common camera trap challenges such as limited training data, strong class imbalance, and heterogeneous backgrounds across camera locations. The authors conduct an ablation study evaluating preprocessing decisions including non-random data splitting, input image resolution, resampling strategies, and data augmentation, followed by a comparison between a classification model (MobileNetV2) and an object detection model (YOLOv5x6 initialized with MegaDetector weights). Experiments are conducted on a dataset of over 250,000 daytime images from 23 camera traps in Switzerland, with only about 1.3\% containing the target species. Results show that the object detection architecture generally outperforms the classification model and achieves higher balanced accuracy and F1-scores, although predictive precision remains low due to extreme class imbalance. The study highlights the importance of transfer learning, oversampling with augmentation, and careful interpretation of performance metrics when applying computer vision models to wildlife monitoring tasks \citep{burkard2024automated}.

\section{Object Detection and Hierarchical Recognition Pipelines}

A related research direction emphasizes object detection as a basis for wildlife recognition, particularly in complex backgrounds or when multiple animals may appear in a scene. \citep{tan2022animal} investigate the application of deep learning–based object detection for automated wildlife recognition in camera trap data, addressing the challenge of efficiently processing large volumes of images collected in ecological monitoring projects. The authors construct the Northeast Tiger and Leopard National Park (NTLNP) dataset, consisting of 25,657 annotated images from infrared camera traps collected between 2014 and 2020, covering 17 species including both wild and domestic animals. Using this dataset, the study compares the performance of three mainstream object detection architectures representing different design paradigms: the anchor-based one-stage YOLOv5 models, the anchor-free one-stage FCOS models with ResNet backbones, and the anchor-based two-stage Cascade R-CNN with HRNet32. Experiments evaluate detection accuracy under different training strategies, including separate and joint training of daytime and nighttime images. Results show that all models achieve high detection performance, with mean average precision above 97.9\% at IoU 0.5, while YOLOv5m achieves the highest overall accuracy. The authors also demonstrate the application of the trained models to camera trap video classification. However, the study notes limitations related to image quality, background similarity, and species with small or occluded appearances, which can reduce detection accuracy in practical monitoring scenarios \citep{tan2022animal}.

Beyond direct end-to-end recognition, some studies propose multi-stage pipelines to address confusion among visually similar species and background-induced errors. \citep{mulero2025addressing} address key challenges in automated animal detection from camera trap images, including class imbalance, confusion between visually similar species, and the influence of background context on classification. To tackle these issues, the authors propose a two-stage deep learning framework that combines a global detection model with specialized expert models. The approach first employs a global model based on transfer learning from MegaDetectorV5 (YOLOv5) to classify images into broad groups of animals. Subsequently, hierarchical agglomerative clustering is used to group species according to visual similarity, and dedicated expert models are trained for each group to perform more precise classification. The method was evaluated on a large dataset comprising approximately 1.3 million images collected from camera traps deployed in Doñana and Sierra de las Nieves National Parks in Spain, covering 24 mammal species. Experiments conducted on an out-of-sample test set of more than 120,000 images showed that the proposed two-stage pipeline achieved an F1-score of 96.2\%, improving upon a single global model (F1-score 0.92). Although effective, the approach relies on expert-defined grouping strategies and extensive labeled data, which may limit scalability to other ecological contexts \citep{mulero2025addressing}.

\section{Multimodal and Metadata-Enhanced Wildlife Recognition}

While most camera trap recognition systems rely exclusively on image content, a growing line of work incorporates contextual information into the recognition process. \citep{liu2024improved} address the limitations of wildlife recognition systems that rely solely on visual features from camera trap images by incorporating contextual temporal information into the classification process. The authors propose a multimodal deep learning framework, termed Temporal-SE-ResNet50, which combines image-based features with temporal metadata derived from camera trap timestamps. The architecture employs an SE-ResNet50 backbone with a squeeze-and-excitation attention mechanism to extract visual features from images, while temporal information such as date and time is encoded using sine–cosine cyclical transformations and processed through a residual multilayer perceptron network. These visual and temporal representations are then fused using a dynamic MLP module to produce the final wildlife classification. Experiments are conducted primarily on the Camdeboo camera trap dataset, containing 9,859 images of 15 animal species, with additional evaluation on subsets of the Snapshot Serengeti and Snapshot Mountain Zebra datasets. The proposed model achieves an accuracy of 93.10\% on the Camdeboo dataset, outperforming several baseline convolutional architectures including ResNet50, VGG19, MobileNetV3, and ConvNeXt. However, improvements vary across species, and performance decreases for certain classes, indicating that the benefits of temporal information may depend on species-specific activity patterns and dataset characteristics \citep{liu2024improved}.

Another recent direction seeks to reduce dependence on large annotated datasets through multimodal foundation models and zero-shot inference. \citep{fabian2023multimodal} investigate the problem of animal species recognition in camera trap images without relying on large labeled training datasets. To address the limitations of supervised approaches, which require extensive expert annotations and often struggle to generalize across regions, the authors propose WildMatch, a zero-shot classification framework based on multimodal foundation models. The method uses a vision--language model to generate detailed textual descriptions of animals in camera trap images, which are then matched against an external knowledge base of species descriptions derived from sources such as Wikipedia. To improve caption quality, the authors introduce an instruction-tuning strategy that combines a small set of human-annotated descriptions with automatically generated captions augmented with expert knowledge. The framework also employs hierarchical classification and self-consistency through multiple caption generations to improve prediction reliability. Experiments on a camera trap dataset collected in the Magdalena Medio region of Colombia show that the proposed approach substantially improves zero-shot performance compared with baseline multimodal models and CLIP-based methods. However, the method requires large multimodal models and incurs higher computational cost than conventional supervised classifiers \citep{fabian2023multimodal}.

\section{Beyond Species Classification: Marked Individuals and Operational Workflows}

Most automated camera trap studies focus on species-level recognition, but some tasks require detecting marked individuals for resighting and demographic analyses. Santangeli et al. address the challenge of efficiently processing large volumes of camera trap images by developing a semi-automated system for detecting individually tagged animals using deep learning. While most prior approaches in camera trap analysis focus on species classification, this study targets the identification of animals marked with patagial wing tags to facilitate resighting and demographic studies. The authors implement a convolutional neural network--based object detection model using the YOLOv3 architecture to identify and localize visible tags in images of Lappet-faced vultures (Torgos tracheliotus) collected from camera traps deployed at water points in Namibia. The training dataset includes 1,379 images containing tagged birds and 817 images without tags, with manual bounding box annotations used to supervise the detection task. The final model achieved strong performance, with a mean average precision of 92.3\% and an optimal F1 score of 85.8\% at a confidence threshold of 0.4. The system correctly classified 88.9\% of test images and substantially reduced manual processing time, potentially saving 11--24 working days per camera per year. However, the approach remains semi-automated because manual inspection is still required to read tag codes, and the model is tailored to a specific tag type and species context \citep{santangeli2022semi}.

In parallel with algorithmic advances, some studies examine how computer vision systems can be integrated into end-to-end monitoring infrastructures. \citep{smith2024man} examine the operational benefits of integrating artificial intelligence--based image classification with connected camera trap systems for wildlife monitoring. The study addresses the logistical and financial challenges associated with manual retrieval and analysis of large volumes of camera trap data. The authors evaluate an AI-based classification platform (eVorta), which uses a convolutional neural network detector designed for fixed-focus camera trap images, and combine it with 4G-connected camera traps that automatically transmit images to a cloud-based processing system. A cost--benefit decision tool is introduced to compare three monitoring scenarios: manual image retrieval and classification, manual retrieval with AI-assisted classification, and fully automated workflows using 4G connectivity and AI processing. The framework is evaluated using a feral cat eradication monitoring program on Kangaroo Island, Australia, involving a network of 200 cameras and a dataset of more than 100,000 images. Results indicate that AI-assisted processing improves classification accuracy relative to human observers and significantly reduces monitoring costs and travel-related carbon emissions when combined with real-time image transmission. However, the approach depends on reliable network connectivity and may require region-specific model training to ensure transferability across monitoring sites \citep{smith2024man}.

\section{Computer Vision for Small Fauna and Insect Monitoring}

Although much of the camera trap literature focuses on vertebrate monitoring, related computer vision methods have also been developed for insects and other small taxa. \citep{bjerge2023accurate} investigate the use of deep learning for automated detection and classification of insects in camera-based ecological monitoring systems. The study addresses the challenge of identifying small and visually diverse insects in complex natural backgrounds, which complicates the extraction of ecological information from image-based monitoring systems. The authors develop a computer vision pipeline based on YOLO object detection models and introduce a large annotated dataset derived from time-lapse camera recordings. The dataset contains 29,960 annotated insect instances across nine taxa, extracted from more than two million images captured by ten cameras mounted above flowering plants during a summer field deployment. The annotated data were created using a multi-stage process that combined automated detection, citizen science screening, and expert validation. Experimental evaluation shows that a YOLOv5 model achieves an average precision of 92.7\% and recall of 93.8\% for detection and classification across the nine taxa. The model can also detect unfamiliar insect species and assign them to related classes. However, the approach is limited by the restricted number of species in the training data and by the similarity of backgrounds in the recorded images, which may affect generalization to other ecological contexts \citep{bjerge2023accurate}.

Similarly, \citep{biswas2024utilizing} investigate the application of computer vision and deep learning for automated monitoring of insects using camera trap systems. The study addresses the limitations of traditional insect monitoring methods, which typically require manual collection and identification of specimens and provide limited temporal resolution. To address this problem, the authors design an automated monitoring system based on a camera-equipped moth trap integrated with a lightweight deep learning pipeline. The system includes a four-tier architecture—perception, transport, processing, and application layers—to support image acquisition, transmission, analysis, and visualization. A computer vision method called Moth Categorization and Counting (MCC) is introduced to detect, track, count, and classify insects in captured images. The classification component relies on a custom convolutional neural network trained on approximately 2,000 labeled images covering eight insect categories, with additional images used for testing and evaluation. The system was deployed in real-world conditions, capturing approximately 250,000 images over 48 nights. Experimental results show classification accuracies around 94--96\% for several insect categories and demonstrate near real-time inference. However, the method is limited by the relatively small labeled dataset and by the restricted number of insect categories considered in the experiments \citep{biswas2024utilizing}.

\section{Final remarks}

The literature reviewed in this chapter shows that camera trap analysis has evolved from manual interpretation toward increasingly sophisticated computer vision pipelines. Current research spans supervised image classification, object detection, hierarchical and two-stage recognition frameworks, metadata-enhanced multimodal models, zero-shot vision--language approaches, and operational systems that integrate automated inference with connected sensing infrastructures. Across these directions, transfer learning, object localization, attention mechanisms, expert specialization, and contextual fusion emerge as recurring strategies for improving recognition performance under ecologically realistic conditions.

Despite these advances, several limitations remain consistent across the literature. Many methods depend on large volumes of labeled data, yet available datasets are often taxonomically narrow, geographically constrained, class-imbalanced, and affected by domain shift across sites, seasons, lighting conditions, and camera configurations. Performance also tends to degrade for visually similar species, small or occluded animals, rare classes, and low-quality images. Multimodal and zero-shot approaches partly address annotation scarcity, but they introduce higher computational cost and do not fully resolve variability across ecological contexts. In addition, most studies emphasize species-level classification, with more limited attention to fine-grained recognition, uncertainty handling, and robust generalization under field deployment conditions.

These gaps directly motivate the research developed in this thesis. In particular, they indicate the need for methods that are more robust to data scarcity and distribution shift, better able to exploit complementary sources of information, and more suitable for fine-grained wildlife monitoring scenarios in ecologically diverse camera trap settings.
